\documentclass[11pt]{article}
\usepackage{multicol} % Required package
\usepackage{lipsum}   % Package for dummy text
\usepackage{amsmath,amsthm,amscd,amssymb,mathrsfs,tikz}
% \usepackage{verbatim}
\usepackage{latexsym,epsf,epsfig}
\usepackage{sectsty}
\usepackage{hyperref}
\usepackage{graphicx}
\usepackage{booktabs, multirow} % for borders and merged ranges
\usepackage{soul}% for underlines
\usepackage{xcolor,colortbl} % for cell colors
\usepackage{changepage,threeparttable} % for wide tables
\usepackage[left=1.2in, right=1.2in, top=1in, bottom=1in]{geometry} %The above set the parameters adjust the margins. There are many ways to do this, and frankly, what is above is not the best. However, it will work the time being.

% \usepackage{csquotes}
\usepackage{tcolorbox}
% Define a custom box called "blockquote"
\newtcolorbox{blockquote}{
    colback=white,    % Background color
    colframe=gray,    % Line color
    boxrule=0pt,      % Width of borders
    leftrule=3pt,     % Width of the left line
    arc=0pt,          % Square corners
    outer arc=0pt,
    left=10pt,        % Space between line and text
    top=0pt,
    bottom=0pt,
    right=0pt
}
\newcommand{\ds}{\displaystyle}
\newcommand{\sU}{\mathscr U}
% Theorem
\theoremstyle{plain}
\newtheorem{theorem}{Theorem}
% Margins
\topmargin=-0.45in
\evensidemargin=0in
\oddsidemargin=0in
\textwidth=6.5in
\textheight=9.0in
\headsep=0.25in

\title{ MATH 301: Homework 2}
\author{ Aren Vista }
\date{\today}

\begin{document}
\maketitle
\pagebreak
\begin{center}
	\section*{Question 1: Section 2.1-1}
	If $a,b \in \mathbb{R},$ prove the following:  \\
\end{center}
\subsection*{\textbf{Part (a):} \textit{If $a+b=0$, then $b=-a$}}
\begin{blockquote}
	\textbf{Recall Properties of $\mathbb{R}$ Addition:} \\
	\textit{(Commutativity:)} $a+b = b+a$ \\
	\textit{(Associativity:)} $(a+b)+c = a+(b+c)$ \\
	\textit{(Zero Property:)} There exists a $0 \in \mathbb{R}$ such that $a+0=a$ \\
	\textit{(Existence of inverse:)} $\forall ~a \in \mathbb{R}, \exists ~-a \in \mathbb{R}$ such that $a+(-a)=0$
\end{blockquote}
Let $a+b=0$. Observe the following:
\[
	\begin{array}{cc}
		\exists ~a,b \in \mathbb{R}, b \neq -a \ni & \text{By proposition}                    \\
		a+b=0                                      & \text{By proposition}                    \\
		a+b+(-a) = 0+(-a)                          & \text{By existance of inverses}          \\
		b+(a+(-a)) = 0+(-a)                        & \text{By associativity}                  \\
		b+(0) = 0+(-a)                             & \text{By inverse property}               \\
		b = -a                                     & \text{By definition of identity element} \\
	\end{array}
\]
\begin{center}
	Clearly we have shown given $a+b=0$, then $b=-a$ holds.
\end{center}
\subsection*{\textbf{Part (c):} \textit{$(-1)a=-a$}}
\begin{blockquote}
	\textbf{Recall Properties of $\mathbb{R}$ Multiplication} \\
	\textit{(Commutativity:)} $a \cdot b = b \cdot a$ \\
	\textit{(Associativity:)} $(a \cdot b) \cdot c = a \cdot (b \cdot c)$ \\
	\textit{(Identity:)} $\exists ~ 1 \in \mathbb{R}$ such that $a \cdot 1 = a$ \\
	\textit{(Existence of inverse:)} $\forall ~a \neq 0, \exists ~\frac{1}{a} \in \mathbb{R}$ such that $a \cdot \frac{1}{a} = 1$ \\
\end{blockquote}
Consider the following:
\[
	\begin{array}{cc}
		a + (-1)a =  a + (-1)a        & \text{ Proposition x=x}              \\
		a + (-1)a = 1 \cdot a + (-1)a & \text{ Multiplicative Identity}      \\
		a + (-1)a = a(1+(-1))         & \text{ Distributive Property}        \\
		a + (-1)a = a(0)              & \text{ Inverse Property of Addition} \\
		a + (-1)a = 0                 & \text{ Zero Multiplicative Property} \\
	\end{array}
\]
\begin{center}
	Thus we can say $a + (-1)a = 0$. \\
	Recall by \textit{Part (a)} "If $a+b=0$, then $b=-a$". \\
	Therefore $(-1)a=-a$ must hold. \\
\end{center}
\subsection*{\textbf{Part (d):} \textit{$-(-a) = a$}}
Consider the following:
\[
	\begin{array}{cc}
		(-a)+a=0           & \text{ Definition of Additive Inverse}     \\
		(-1)((-a)+a)=(-1)0 & \text{ Multiply both sides by $(-1)$}      \\
		-(-a)+(-a)=(-1)0   & \text{ Distributive Property and Part (b)} \\
		-(-a)+(-a)=0       & \text{ Multiplicative Property of Zero}    \\
		-(-a)+(-a)+a=0+a   & \text{ Add $a$ to both sides}              \\
		-(-a)+((-a)+a)=a   & \text{ Communitive Property}               \\
		-(-a)+(0)=a        & \text{ Inverse Property of Addition}       \\
		-(-a)=a            & \text{ Conclusion}                         \\
	\end{array}
\]
\begin{center}
	Thus we have shown that $-(-a)=a$.
\end{center}
\subsection*{\textbf{Part (e):} \textit{$(-1)(-1)=1$}}
Consider $-(-a) = a$, we stated $a \in \mathbb{R}$ as an arbitrary element.
\[
	\begin{array}{cc}
		-(-a)=a            & \text{ Part (d)}                                              \\
		-(-1(a)) = 1(a)    & \text{Identity Property of Multiplication}                    \\
		(-1)(-1(a)) = 1(a) & \text{Part (c)}                                               \\
		(-1)(-1(1)) = 1(1) & \text{Let $a=1$, $a$ is an arbitrary element of $\mathbb{R}$} \\
		(-1)(-1) = 1       & \text{Identity Property of Multiplication}                    \\
	\end{array}
\]
\begin{center}
	Thus we have shown $(-1)(-1) = 1$
\end{center}
\begin{center}
	\section*{Question 2: Section 2.1-2}
	If $a,b \in \mathbb{R},$ prove the following:  \\
\end{center}
\subsection*{\textbf{Part (a):} $-(a+b) = (-a) + (-b)$}
Consider the following:
\[
	\begin{array}{cc}
		-(a+b) = (-1)(a+b)    & \text{By q1.c}                  \\
		-(a+b) = (-1)a+ (-1)b & \text{By Distributive Property} \\
		-(a+b) = -a+ -b       & \text{By q1.c}                  \\
	\end{array}
\]
\subsection*{\textbf{Part (b):} $(-a) \cdot (-b) = a \cdot b$}
Consider the following:
\[
	\begin{array}{cc}
		(-a) \cdot (-b) = (-a) \cdot (-b)                 & \text{Proposition x=x}                        \\
		(-a) \cdot (-b) = (-1) \cdot a \cdot (-1)\cdot b  & \text{By q1.c}                                \\
		(-a) \cdot (-b) = (-1) \cdot (-1) \cdot a \cdot b & \text{By Associativity}                       \\
		(-a) \cdot (-b) = (1) \cdot a \cdot b             & \text{By q1.e}                                \\
		(-a) \cdot (-b) =  a \cdot b                      & \text{By Identity Property of Multiplication} \\
	\end{array}
\]
\subsection*{\textbf{Part (c):} $\frac{1}{(-a)}= -(\frac{1}{a})$}
Consider the following:
\[
	\begin{array}{cc}
		(\frac{1}{-a})= -(\frac{1}{a})      & \text{By Proposition}                      \\\\
		(\frac{1}{-a})= (-1)(\frac{1}{a})   & \text{By q1.c}                             \\\\
		(\frac{1}{-a})a= (-1)(\frac{1}{a})a & \text{Multiply both sides by $b$}          \\\\
		(\frac{1}{-a})a= (-1)1              & \text{Property of Multiplicative Inverse}  \\\\
		(\frac{1}{-a})a= (-1)               & \text{Property of Multiplicative identity} \\\\
		(\frac{1}{-a})(-1)a= (-1)(-1)       & \text{Multiply both sides by (-1)}         \\\\
		(\frac{1}{-a})(-a)= (-1)(-1)        & \text{By q1.c}                             \\\\
		1 = (-1)(-1)                        & \text{Property of Multiplicative Inverse}  \\\\
		1 = 1                               & \text{By q1.e}                             \\\\
	\end{array}
\]
\begin{center}
	\section*{Question 3: Section 2.1-10}
\end{center}
\subsection*{\textbf{Part (a):} If $a<b$ and $c \leq d$, prove that $(a+c)<(b+d)$}
Suppose $a<b$ and $c \leq d$. Consider the following:
\begin{blockquote}
	\textbf{Recall: } \textit{(Theorem) Rules of Inequality} \\
	$\text{Let } a,b,c \in \mathbb{R}$. The following statments must hold:
	\begin{multicols}{2}
		\begin{itemize}
			\item $(a>b) \land (b>c) \implies a>c$
			\item $(a>b) \implies (a+c) > (a+b)$
			\item $(a>b) \land (c>0) \implies ca > cb$
			\item $(a>b) \land (c<0) \implies ca < cb$
		\end{itemize}
	\end{multicols}
\end{blockquote} \noindent
Applying the theorem above, consider the following:
\[
	\begin{aligned}
		 & a<b                 \\
		 & a+c < b+c           \\
		 & b+c \leq  b+d       \\
		 & a+c < b+c \leq  b+d \\
		 & a+c <  b+d          \\
	\end{aligned}
\]

\subsection*{\textbf{Part (b):} If $0<a<b$ and $0 \leq c \leq d$, prove that $0 \leq ac \leq bd$}
Suppose $0<a<b$ and $0 \leq c \leq d$
\begin{blockquote}
	\textbf{Recall: } \textit{(Theorem) Rules of Inequality} \\
	$\text{Let } a,b,c \in \mathbb{R}$. The following statments must hold:
	\begin{multicols}{2}
		\begin{itemize}
			\item $(a>b) \land (b>c) \implies a>c$
			\item $(a>b) \implies (a+c) > (a+b)$
			\item $(a>b) \land (c>0) \implies ca > cb$
			\item $(a>b) \land (c<0) \implies ca < cb$
		\end{itemize}
	\end{multicols}
	Ordering
	If $a-b \in P$, we write $a>b$ (or $b<a$). \\
	If $a-b \in P \cup \{ 0 \}$, we write $a \geq b$ (or $b \leq a$). \\

	For any two real numbers $a,b$ exactly one of the following holds:
	$$a>b, a<b, a=b$$
\end{blockquote} \noindent
\begin{blockquote}
	\textbf{Recall: } \textit{ Properties of Positive Sets} \\
	Let $P \neq \emptyset, P \subseteq \mathbb{R}$ be called the set of positive real numbers. \\
	1.  If $a,b \in P$ then $a+b \in P$.\\
	2.  If $a,b \in P$ then $a \cdot b \in P$.\\
	3.  \textit{Trichotomy Property:} If $a \in \mathbb{R}$, then exactly one of the following must be true:
	$$a \in P, ~~a = 0, ~-a \in P$$
\end{blockquote} \noindent
Applying the theorems above, observe the following:
\[
	(0 < a) \land (0 \leq c) \implies (0 \leq ac)
\]
Thus the following statements hold:
\[
	\begin{aligned}
		 & 0 \leq c \leq d            \\
		 & 0 \leq bc \leq bd          \\
		 & 0 < a < b                  \\
		 & 0 \leq ac \leq  bc \leq db \\
		 & 0 \leq ac \leq db          \\
	\end{aligned}
\]
\begin{center}
	\section*{Question 4: Section 2.1-11}
\end{center}
\subsection*{Part (a): If $(a>0)$, then $(1/a>0)$ and $\frac{1}{1/a} = a$}
Suppose $a>0$ and $\frac{1}{a} \not > 0$. Consider the following:
\[
	\begin{aligned}
		\frac{1}{a}           & \not > 0         \\
		(\frac{1}{a}) \cdot a & \not > 0 \cdot a \\
		(\frac{1}{a}) \cdot a & \not > 0         \\
		1                     & \not > 0         \\
	\end{aligned}
\]
\begin{center}
	Clearly $1>0$ thus this is a contradiction. \\
	Thus we must accept the original proposition: $a>0 \implies \frac{1}{a}>0$.
\end{center}
Prove the following statements $\frac{1}{1/a}=a$:
\[
	\begin{aligned}
		\frac{1}{1/a}                & =a              \\
		(\frac{1}{1/a})(\frac{1}{a}) & =a(\frac{1}{a}) \\
		1                            & = 1
	\end{aligned}
\]
\begin{center}
	Clearly $1=1$ thus proven directly.
\end{center}
Now we have proven if $(a>0)$, then $(1/a>0)$ and $\frac{1}{1/a} = a$.
\subsection*{Part (b): If $(a<b)$, then $a<\frac{1}{2}(a+b)<b$}
Consider the following:
\[
	\begin{aligned}
		a                           & =  a                           \\
		\frac{1}{2}a                & =  \frac{1}{2}a                \\
		\frac{1}{2}a + \frac{1}{2}a & =  \frac{1}{2}a + \frac{1}{2}a \\
		\frac{1}{2}(a + a)          & =  \frac{1}{2}(a + a)          \\
		\frac{1}{2}(2a)             & =  \frac{1}{2}(a + a)          \\
		\frac{1}{2}2(a)             & =  \frac{1}{2}(a + a)          \\
		a                           & =  \frac{1}{2}(a + a)          \\
		a <                         & \frac{1}{2}(a + b) < b         \\
	\end{aligned}
\]
\begin{center}
	\section*{Question 5: Sec 2.1-13}
\end{center}
If $a,b \in \mathbb{R}$, show that $a^2+b^2=0$ iff $a=0$ and $b=0$.  \\
\subsection*{Case 1: $\rightarrow$}
\textbf{Prove:} If $a,b \in \mathbb{R}$, show that $a^2+b^2=0$ implies $a=0$ and $b=0$.  \\
Consider the contrapositive statement, $a \neq 0$ or $b \neq 0$ implies $a^2 + b^2 \neq 0$:\\
Observe the following: \\

If $a \in \mathbb{P}$
\[
	a \in \mathbb{P} \implies a^2 = a \cdot a  \in \mathbb{P} \implies a^2 > 0
\]

If $-a \in \mathbb{P}$: (Consider $\forall ~n \in \mathbb{R}, (n-n) \cdot (-n) = 0$) \\
\[
	\begin{aligned}
		n \cdot (-n) + (-n)(-n) = 0 \\
		(-n) + (-n)(-n) = 0         \\
		\implies (-n)(-n) = n
	\end{aligned}
\]
\begin{center}
	Therefore, $-a^2 \in \mathbb{P}$. The same logic can be applied to $b$. Observe If $a,b \in \mathbb{P} \implies a+b \in \mathbb{P}$. \\
	Therefore we have proven the contrapositive statement if $a \neq 0$ or $b \neq 0$ implies $a^2 + b^2 \neq 0$
\end{center}
\subsection*{Case 2: $\leftarrow$}
\textbf{Prove:} If $a,b \in \mathbb{R}$, show that $a=0$ and $b=0$ implies $a^2+b^2=0$. \\
Consider the following:
\[
	0 \cdot 0 = 0 \implies \text{If } x=0, \text{ then } x^2 = x \cdot x = 0 \cdot 0 = 0
\]
As $a=b=0$ then naturally if $a=0$ and $b=0$ implies $a^2+b^2=0$
\begin{center}
	\section*{Question 6: Section 2.1-14}
\end{center}
If $0 \leq a<b$, show that $a^2 \leq ab < b^2$. Show by example that it does \textbf{not} follow that $a^2 <ab < b^2$. \\
Let $a=0, b=1$. Thus $0 \leq 0 < 1$ satisfy $a^2 \leq ab < b^2$. Observe $a^2 <ab < b^2$ does not hold as $0 < 0 < 1$ is false as $ 0 \not < 0$

% \begin{center}
\section*{Question 7: Section 2.1-20}
% \end{center}
\begin{center}
	\textbf{Part a: } If $0<c<1$, show that $0<c^2<c<1$. \\
\end{center}
We can define $1<c$ as $1<1-\epsilon$ where $c=1-\epsilon = 1 - \frac{1}{n}, n \in \mathbb{N}$. \\
Observe the following: \\
\[
	\begin{aligned}
		1 < (1-\epsilon)                                       \\
		1^2 < (1-\epsilon)^2                                   \\
		1^2 < (1-\frac{1}{n})^2                                \\
		(1-\frac{1}{n})^2  = 1 - \frac{2}{n} + \frac{1}{n^2}   \\
		\text{Observe } \frac{1}{n} > \frac{1}{n^2} \implies   \\
		1 > (1-\frac{1}{n}) > (1- \frac{2}{n} + \frac{1}{n^2}) \\
		\text{Thus we can conclude: } 1 > 1-\epsilon > (1-\epsilon)^2
	\end{aligned}
\]

\begin{center}
	\textbf{Part b: } If $1<c$, show that $1<c<c^2$.
\end{center}
We can define $1<c$ as $1<1+\epsilon$ where $c=1+ \epsilon = 1 + \frac{1}{n}, n \in \mathbb{N}$. \\
Observe the following: \\
\[
	\begin{aligned}
		1 < (1+\epsilon)                      \\
		1^2 < (1+\epsilon)^2                  \\
		(1+\epsilon)^2 = 1 + 2 \epsilon + 1   \\
		1 < 1 + \epsilon < 1 + 2 \epsilon + 1 \\
		1 < c < c^2
	\end{aligned}
\]
As we defined $\epsilon > 0$.
\[
	\text{The statement } 1 < c < c^2 \text{ holds}
\]
% \begin{center}
\section*{Question 8: Section 2.1-23}
% \end{center}
If $a>0$, $b>0$, and $n \in \mathbb{N}$, show that $a<b$ iff $a^n < b^n$ (induct.) \\ \\
$\forall ~a,b \in \mathbb{R}, a>0, b>0$ and $n \in \mathbb{N}$. \\
Consider $P(n): a^n < b^n \iff a < b$ \\
Observe $P(1)$:
\[
	\begin{aligned}
		a^1 < b^1 \implies a < b \\
		a < b \implies  a^1 < b^1
	\end{aligned}
\]
Suppose $P(k), k \in \mathbb{N}$:
\[
	a^k < b^k \implies  a < b
\]
Observe the following:
\[
	\begin{aligned}
		a^k < b^k \implies  a < b      \\
		a^k < b^k                 \\
		a^k \cdot ab < b^k \cdot ab \\
        a^{k+1}b < b^{k+1}a \\
	\end{aligned}
\]
% \begin{center}
	\section*{Question 9: Section 2.1-25}
% \end{center}
%
% \begin{center}
% 	\section*{Question 9: Section 2.2-1}
% \end{center}
%
% \begin{center}
% 	\section*{Question 9: Section 2.2-2}
% \end{center}
%
% \begin{center}
% 	\section*{Question 9: Section 2.2-10}
% \end{center}
%
% \begin{center}
% 	\section*{Question 9: Section 2.2-12}
% \end{center}
%
% \begin{center}
% 	\section*{Question 9: Section 2.2-17}
% \end{center}
%

\end{document}
